\documentclass[12pt]{article}
\usepackage{amsmath, amssymb, amsthm}
\usepackage[margin=1in]{geometry}

% --- Theorem environments ---
\newtheorem{theorem}{Theorem}
\newtheorem{definition}{Definition}

\title{\textbf{Divisibility of the Sum of Three Consecutive Integers}}
\author{0xHaru}
\date{\today}

\begin{document}
\maketitle

\begin{definition}[Divisibility]
    We say $a \mid b$ (\emph{``$a$ divides $b$''}) if there exists some $k \in \mathbb{Z}$ such that
    \[
        b = k \cdot a.
    \]
\end{definition}

\begin{theorem}
    The sum of any three consecutive integers is divisible by $3$.
\end{theorem}

\begin{proof}
    Let $n \in \mathbb{Z}$, and let $n,\ n+1,\ n+2$ be three consecutive integers. Their sum is
    \[
        n + (n+1) + (n+2) = 3n + 3 = 3(n+1).
    \]
    Since $n + 1 \in \mathbb{Z}$, taking $k = n + 1$ in Definition~1 yields $3 \mid 3(n+1)$, and hence the sum is divisible by $3$.
\end{proof}

\end{document}
